\setscript[hanzi]
\setupalign[hanging, hz]

\usemodule[zhfontsext]
\usetypescript[cnfontb]
\setupbodyfont[cnfontb,rm,12pt] % 小四号 = 12pt
\mainlanguage[cn]

\setupinterlinespace[line=3.66ex] %line=3.05ex相当于WORD 中的1.0倍行距, 3.66ex相当于WORD 中的1.2倍行距

\setuppapersize[A4]
\setuplayout[width=fit,height=fit,backspace=10mm,cutspace=10mm,topspace=10mm,margin=10mm,header=10mm,footer=10mm]
\setuppagenumbering[location={header,right},style=\bfb]

\setupinteraction[state=start]

% % % ================== 脚注 尾注 ==================

\startsetups footnoteEnv
\switchtobodyfont[10pt]
\setscript[hanzi]
\stopsetups

\startsetups endnoteEnv
\switchtobodyfont[10pt]
\setscript[hanzi]
\stopsetups



\starttext

% % % ------------------ 段落中的脚注尾注 ------------------

% % % ================== 代码示例 - 1


\setupnote[footnote][setups=footnoteEnv]
\setupnote[endnote][setups=endnoteEnv]


\midaligned{
三字经\footnote[title]{中国传统的儿童启蒙读物，用简洁通俗的白话讲出了亘古不变的哲理，脍炙人口、广为流传.}\\
人之初，性本善。性相近，习相远\note[sentence1]。\quad
苟不教，性乃迁。教之道，贵以专。\\
昔孟母，择邻处。子不学，断机杼\endnote[sentence3]{战国时，孟子的母亲为了让孟子有个好的学习环境搬了三次家。有一次孟子厌学，孟母就割断织布机上的梭子来教育他，做事情要持之以恒，不能半途而废，否则会前功尽弃。}。\quad
窦燕山，有义方。教五子，名俱扬。\\
养不教，父之过。教不严，师之惰。\quad
子不学，非所宜。幼不学，老何为。\\
玉不琢，不成器。人不学，不知义。\quad
为人子，方少时。亲师友，习礼仪。\\
香九龄，能温席。孝于亲，所当执。\quad
融四岁，能让梨，悌于长，宜先知。\\
首孝悌，次见闻。知某数，识某文。\quad
一而十，十而百。百而千，千而万。\\
三才者，天地人。三光者，日月星。\quad
三纲者，君臣义。父子亲，夫妇顺。\\
}
\footnotetext[sentence1]{人刚出生的时候，禀性是善良。每个人的禀性本来是很接近的，只是由于后天的生活环境和学习环境不同，差异就越来越大}

在\note[sentence3]中，提到了孟母三迁的故事。
\placenotes[endnote]


\stoptext